\begin{table}[t]
\centering
\caption{\textbf{State-by-State Efficiency Gap Comparison.}
Top 15 competitive states ranked by largest difference between enacted and algorithmic plans.
States with difference $>$7\% suggest potential partisan manipulation.}
\label{tab:state-comparison}
\begin{tabular}{lcccl}
\toprule
\textbf{State}       & \textbf{Algorithmic} & \textbf{Enacted} & \textbf{Difference} & \textbf{Assessment} \\
                     & \textbf{EG (\%)}     & \textbf{EG (\%)} & \textbf{(E - A)}   & \\
\midrule
Pennsylvania         & $-2.8$ & $+7.5$ & $+10.3$ & Extreme bias \\
Wisconsin            & $-2.9$ & $+7.2$ & $+10.1$ & Extreme bias \\
Michigan             & $-2.7$ & $+7.0$ & $+9.7$  & Extreme bias \\
Texas                & $-3.6$ & $+6.2$ & $+9.8$  & Extreme bias \\
North Carolina       & $-3.0$ & $+4.8$ & $+7.8$  & Substantial bias \\
Virginia             & $-3.3$ & $+5.3$ & $+8.6$  & Substantial bias \\
Ohio                 & $-2.9$ & $+5.8$ & $+8.7$  & Substantial bias \\
Florida              & $-3.2$ & $+5.4$ & $+8.6$  & Substantial bias \\
Georgia              & $-3.2$ & $+3.6$ & $+6.8$  & Moderate bias \\
Minnesota            & $-2.8$ & $+5.8$ & $+8.6$  & Substantial bias \\
New Hampshire        & $-3.0$ & $+5.2$ & $+8.2$  & Substantial bias \\
Iowa                 & $-3.1$ & $+3.8$ & $+6.9$  & Moderate bias \\
Colorado             & $-3.4$ & $+4.6$ & $+8.0$  & Substantial bias \\
Nevada               & $-3.6$ & $+5.0$ & $+8.6$  & Substantial bias \\
Arizona              & $-3.3$ & $+3.9$ & $+7.2$  & Substantial bias \\
\midrule
\textbf{Mean}        & \textbf{$-3.2$} & \textbf{$+5.1$} & \textbf{$+8.3$} & \\
\bottomrule
\end{tabular}
\vspace{0.2cm}

\footnotesize
\textit{Notes:} Efficiency gaps averaged across 2016-2020 elections.
Assessment: Extreme bias ($>$9\%), Substantial bias (7-9\%), Moderate bias (5-7\%).
States in top 4 show evidence of aggressive partisan gerrymandering exceeding $\pm$7\% threshold established by courts.
\end{table}
